The following bounds apply to the same completed displacement on its
actual connected through interval. The appendices prove the transition
and section estimates used in each region, with the one-sided endpoint
estimates and continuations needed to count the full physical carrier.

\subsection{The signed high-height interval}

\begin{corollary}[Signed high-height interval]
\label{rew:cor-high-derivative-signed}
There exist \(K,z_*,\gamma_*,C_H>0\) such that, for all sufficiently
small raw coefficients and either sign of \(\alpha_0\), the complete
transition on the actual fixed sections satisfies \(0<\Pi'\le C_H\)
throughout the interval of through trajectories
\begin{equation}\label{ann:high-source-interval}
 K(|\alpha_0|+\alpha_1^2)<H:=z(\gamma)<z_*,
 \qquad |\gamma|<\gamma_*.
\end{equation}
After the common choice \(mC_H<1/2\), the complete displacement has
\(\Delta'>1/2\) and hence at most one isolated zero, counted with
multiplicity, on this entire interval and on each of its truncations.
The assertion includes \(\alpha_0=\alpha_1=0\). At a positive
boundary height \(H=K(|\alpha_0|+\alpha_1^2)<z_*\), with the same
coefficient bounds, the derivative bound still holds; this equality
face is assigned to the low-height region.
\end{corollary}
\begin{proof}
Fix the coefficients and a through trajectory in
\eqref{ann:high-source-interval}. Freeze \(r_0=\sqrt H\) and put
\[
 a=\frac{\alpha_0}{H},\qquad
 p=\frac{\alpha_1}{\sqrt H},\qquad s=r_0,
 \qquad v=\gamma+r_0X,\quad z=r_0^2Y.
\]
The reference orbit has \(Y(0)=1\), and
\[
 |a|<K^{-1},\qquad |p|<K^{-1/2}.
\]
Choose \(K\) so large that the closed block with the corresponding
nonstrict bounds lies inside the signed compact strict family of
Proposition~\ref{sharp:upper-high}. Then choose \(z_*\) and
\(\gamma_*\) small enough for its raw-tail estimates. Its complete
derivative bound applies with the frozen scale \(r_0\), including
both signs of \(a\) and \(a=p=0\). The phase-independent constant bounds the derivative of the same
physical map, with \(r_0\) frozen during differentiation.

Scalar first variation makes \(H\) increasing with that input, so
\eqref{ann:high-source-interval} and its truncations are intervals.
Choose \(mC_H<1/2\) last; then \(\Delta'>1/2\) throughout.
The closed block includes the positive equality face, assigned to
low height. For double-zero coefficients every positive height is high.
\end{proof}

\subsection{The compact positive-ratio whole-carrier bound}

\begin{proposition}\label{sharp:upper-annulus}
For each fixed \(0<p_0<P<\infty\), the complete captured carrier
with \(a=-1\) and \(p_0\le p\le P\) has at most two
isolated displacement roots, counted with multiplicity.
\end{proposition}
\begin{proof}
Use the one connected actual receiver and its open preimages from
Lemma~\ref{ann:lem-compact-positive-domains}. Keep the physical
affine input \(x\) throughout. On \(I_E\cup I_M\),
\(\Pi_{xx}<0\), so for the physical multiplier \(m>0\),
\(\Delta''=-m\Pi''>0\). On \(I_H\), the complete derivative
bound of Lemma~\ref{ann:p0-paired-high-derivative} gives
\(\Delta'>1/2\) after the single final choice
\(m<\min\{m_{\rm ad},1/(2C_H)\}\).
Every critical point has \(\Delta''>0\).
Lemma~\ref{sharp:upper-critical-count} gives the whole-carrier
multiplicity bound.
\end{proof}
\subsection{The gamma-dominant and far-root carriers}

\begin{proposition}\label{sharp:upper-gamma}
For a fixed \(\Pi_{\max}<\infty\), the bounded-\(\pi\)
negative-gamma carrier has at most two displacement zeros, and the
bounded-\(\pi\) positive-gamma carrier has at most three. The positive
far-root carrier of \eqref{deg:positive-far-range} has at most two.
All counts hold on the entire captured physical interval.
\end{proposition}
\begin{proof}
Propositions~\ref{deg:prop-negative-gamma} and
\ref{deg:prop-positive-far} give the two-zero counts on their
original connected carriers, with multiplicity, using the covers
\eqref{deg:negative-gamma-cover} and \eqref{deg:far-physical-cover}.
For positive bounded \(\pi\),
Corollary~\ref{deg:cor-positive-quadratic-whole-carrier} gives at most
two zeros of the same physical \(F'\), including the ghost raw tail.
Its paired/high overlap, in \(\eta=gY_0\), is
\begin{equation}\label{sharp:upper-gamma-overlap}
 gK_H(1+\Pi_{\max}^2g)<\eta_{\max}/2,\qquad
 (gk)^2\eta_{\max}<H_{\rm cap}/2.
\end{equation}
The remaining root-free fibres are contractions by
Lemma~\ref{deg:lem-raw-jacobian}. Lemma~\ref{sharp:upper-critical-count}
then gives the three-zero count on the whole original carrier.
\end{proof}

\subsection{The compact cubic block: a common curvature orientation}
\label{sharp:upper-cubic-section}

\begin{proposition}[The complete compact cubic carrier]
\label{sharp:upper-cubic}
Fix finite \(U_0,G_0\). On the coefficient block
\begin{equation}\label{sharp:upper-cubic-block}
 a=-1,\quad p=tU,\quad\gamma=t^2G,\quad s=t^3,
             \qquad0\le U\le U_0,\quad |G|\le G_0,
\end{equation}
the same displacement has at most two critical points and at most three
isolated zeros, counted with multiplicity, on its whole physical carrier
under the common positive unscaled raw cap. All captured target cycles
lie in this carrier after the common coefficient and marked
Hausdorff-collar shrink.
\end{proposition}

\begin{proof}
Put \(\tau=t^2\). Proposition~\ref{deg:prop-compact-cubic} supplies the
actual clipped carrier and its allocation \eqref{deg:compact-physical-cover}.
Its bounded-action and independent root-collar estimates give the same
curvature orientation. The growing range has positive curvature by
\eqref{deg:compact-curvature}, and is identified with the bounded
range by the actual resection
\begin{equation}\label{sharp:upper-cubic-overlap}
 p_\tau(u)=u+O_{C^3}(\tau(1+|\log\tau|)),\quad p_\tau'>0,
                \qquad2M+\beta\le u\le3M-\beta.
\end{equation}
These are estimates for the completed physical map, so
\begin{equation}\label{sharp:upper-cubic-common-crossing}
 \Pi''(x)=0\quad\Longrightarrow\quad\Pi'''(x)<0
\end{equation}
on the connected lower receiver. Real and fold fibres retain their
one-sided through endpoints.

The original ghost tail has the derivative bound
\eqref{deg:compact-raw-tail-derivative}; the core has
\eqref{deg:compact-core-variation}. The construction places
their comparison and high intervals in strict overlap:
\begin{equation}\label{sharp:upper-cubic-height-overlap}
 \eta_{\rm comp}<\tau K_H(1+\tau U^2)<\eta_{\max}/2,
                      \qquad\tau^2\eta_{\max}<h_c/2.
\end{equation}
One final multiplier bound makes \(F'>1/2\) on this complement.
The cover in Proposition~\ref{deg:prop-compact-cubic} includes every
input, retaining the original clips and equality dividers. Its whole-carrier
count, equivalently Lemma~\ref{sharp:upper-curvature-count}, gives
the stated multiplicity bounds. Capture is applied after these analytic
choices, as in Lemma~\ref{lem:loc-capture}.
\end{proof}

\subsection{Degree loss: a common orientation of critical points}
\label{sharp:upper-degree-section}

The compact graph \eqref{deg:degree-graph} and reciprocal graph
\eqref{deg:degree-infinity-graph} share the physical carrier of
Proposition~\ref{p2:prop-unified-carrier-cover}. The reciprocal face
\(j=0\) concerns only the one-sided closure of the estimates.

\begin{proposition}[The complete degree-loss carrier]
\label{sharp:upper-degree}
Fix the compact positive-unit ranges in
Proposition~\ref{deg:prop-degree-drop}, or the ranges of \(d,A_d\) and
\(0<j\le j_0\) in Proposition~\ref{deg:prop-degree-infinity}.
There are independent active coefficient bounds, one fixed raw-height
cap and a final finite-closing multiplier bound for which the complete
displacement has at most one critical point on its connected interval
of actual through inputs in the original nested windows.
Every such point is a simple upward crossing of the same physical
\(F'\). The displacement therefore has at most two isolated zeros,
counted with multiplicity, and no identity interval.
All estimates hold at every retained input, including nonclosed
through orbits and positive divider equalities. They are uniform for
physical \(j>0\) as \(j\downarrow0\).
\end{proposition}

\begin{proof}
Write the return in the affine input of Lemma~\ref{lem:closing} as
\begin{equation}\label{sharp:upper-degree-displacement}
 F_\theta(x)=x-f_\theta\bigl(m\mathfrak T_\theta(a_\theta+x)\bigr)
             =x-m\Pi_{\theta,m}(x).
\end{equation}
Here $\mathfrak T_\theta=d_\theta\circ T_\theta^\infty$ and
$f_\theta=C_\theta-C_\theta(0)$.
Proposition~\ref{p2:prop-unified-carrier-cover} gives its actual
clipped carrier. Proposition~\ref{p2:prop-physical-degree-count}
retains both gains, output translation and ordinary $C^3$ closing,
and proves
\begin{equation}\label{sharp:upper-degree-critical-orientation}
 F_\theta'(x)=0\ \Longrightarrow\ F_\theta''(x)>0,
 \qquad \#\{x:F_\theta'(x)=0\}\le1.
\end{equation}
Its whole-carrier count gives the remaining conclusions.
\end{proof}

\subsection{The determinant passage}

\begin{proposition}\label{sharp:upper-determinant}
Fix \(P_D<\infty\), the passive compact coefficient block and original
window reserves. For independent small \(r_D>0\), \(|\gamma|\), with
\(\alpha_0=r_D^2\), \(|\alpha_1/r_D|\le P_D\), the displacement on
\eqref{deg:det-actual-receiver} has at most one critical point and at most
two isolated zeros, counted with total multiplicity, after a common
coefficient and final multiplier shrink. The raw cap is fixed, positive
and unscaled. The assertion holds for every original window, speed or
height clip and includes saddle ratio one.
\end{proposition}

\begin{proof}
Use the single affine \(F(x)=x-m\Pi(x)\) of Lemma~\ref{lem:closing}.
Lemma~\ref{deg:lem-determinant-preparation} supplies the full stopped
composition, both exact absolute gains and ordinary \(C^3\) divided
closing. On the strict and high sets in \eqref{deg:det-height-cover},
choose \(m\) last so that their complete derivative bounds give
\(F'>1/2\). Thus every critical point lies in the actual stopped range.

There \eqref{sharp:upper-det-absolute} retains every factor of the same
physical derivative:
\[
 \Pi'(x_b+\sigma_+\delta)=\mathcal A_\theta\delta^q e^{R_\theta(\delta)},
 \quad q=\wp-1,\quad
 |R_\theta|+|\delta R_\theta'|\le C\delta^\alpha(1+|\log\delta|).
\]
The quotient of the two gains is a unit, so \(\mathcal A_\theta\)
is bounded above and below; neither absolute gain needs a lower bound.
Let \(C_0\) bound \(|\log\mathcal A_\theta+R_\theta|\).
At a critical point \(m\Pi'=1\), set \(t=-\log\delta>0\). Then
\[
 -qt\ge L:=\log(1/m)-C_0.
\]
For \(L>0\), this requires \(q<0\) and \(|q|\ge L/t\);
in particular \(\wp=1\) produces no critical point. Moreover
\[
 \frac{|\delta R_\theta'|}{|q|}
 \le\frac C L t e^{-\alpha t}(1+t)\le C_1/L.
\]
Choose \(L>2C_1\). The Euler derivative of the logarithm of the full
physical gain satisfies
\(\delta\partial_\delta\log\Pi'<q/2<0\) at every critical point.
Since \(dx/d\delta=\sigma_+>0\),
\begin{equation}\label{sharp:upper-det-critical}
 F'(x)=0\quad\Longrightarrow\quad F''(x)=-m\Pi''(x)>0.
\end{equation}
The actual stopped/strict and low/high cover
\eqref{deg:det-height-cover} retains every clipped and intervening
input. Lemma~\ref{sharp:upper-critical-count} now gives at most
one critical point and two displacement zeros.
\end{proof}

\subsection{The signed trace passage}

\begin{proposition}\label{sharp:upper-trace}
Fix \(0<p_0<p_1<\infty\), the passive compact block and original
window reserves. For \(\alpha_0=r^2a\), \(\alpha_1=rp\),
\(|a|\le a_*\), \(p_0\le|p|\le p_1\), independent small
\(r>0\), \(|\gamma|\), and a final common multiplier shrink,
the displacement on the original receiver
\eqref{deg:trace-actual-receiver} has at most two isolated zeros for
\(p>0\), and at most one for \(p<0\), counted with total multiplicity.
Its positive unscaled raw cap and all window, speed and height clips
are retained. Both opening signs, zero opening and closed trace endpoints
are included.
\end{proposition}

\begin{proof}
Proposition~\ref{deg:prop-stopped-trace} gives \(F'>1/2\) for
\(p<0\) and \(F'=0\Rightarrow F''>0\) for \(p>0\) on the whole
original clipped receiver. Lemma~\ref{sharp:upper-critical-count}
gives the stated multiplicity bounds.
\end{proof}

\begin{corollary}\label{sharp:upper-local-coefficient-ranges}
The nonnegative-opening sector \(\alpha_0\ge0\) has at most
two captured displacement zeros. There is a fixed \(P_*<\infty\)
such that the negative-opening sector
\(\alpha_0=-s^2<0\), \(p_{\rm src}=\alpha_1/s\ge P_*\)
also has at most two. Both assertions concern the whole physical
interval and count multiplicities.
\end{corollary}

\begin{proof}
First choose the signed trace collar of
Proposition~\ref{sharp:upper-trace}, for example with
\(1/2\le|p|\le2\). Choose \(P_D>2/\sqrt{a_*}\).
If \(\alpha_0>0\) and
\(|\alpha_1|\le P_D\sqrt{\alpha_0}\), use
Proposition~\ref{sharp:upper-determinant} with
\(r_D=\sqrt{\alpha_0}\). On the complement put
\(r=|\alpha_1|\). Then
\[
 p=\sgn\alpha_1,\qquad
 a=\alpha_0/\alpha_1^2<P_D^{-2}<a_*.
\]
This is the exact trace scaling of
Lemma~\ref{deg:lem-positive-atlas}; the full physical domains are
retained in Lemma~\ref{deg:lem-nonnegative-domain}. The same choice treats
\(\alpha_0=0\), \(\alpha_1\ne0\), with \(a=0\).
Each alternative bounds the displacement on its whole coefficient fibre.

The double-zero coefficient face is a contraction by
Corollary~\ref{rew:cor-high-derivative-signed} at every positive
height, without introducing a determinant scale.

Finally suppose \(\alpha_0=-s^2<0\) and
\(p_{\rm src}=\alpha_1/s>0\). With \(r=\alpha_1\), the
trace coefficients are exactly
\[
 p=1,\qquad a=-p_{\rm src}^{-2}.
\]
Choose \(P_*>2/\sqrt{a_*}\) before the subsequent compact
positive-source interval is fixed. Then
\(p_{\rm src}\ge P_*\) belongs to
Proposition~\ref{sharp:upper-trace}. Take the maximum of the
whole-fibre bounds across these coefficient alternatives.
\end{proof}

